\begin{workedexample}[Worked Example: Data rate and decimation]
An SDR digitizes at $f_s = 100\ \text{Msps}$ with $16$-bit samples. Complex
(I and Q) sampling produces a raw rate of
\[ 100\times10^{6}\times16\times2 = 3.2\ \text{Gbit/s}, \]
far too much to process directly. A digital downconverter selects the wanted
channel and decimates by $50$, leaving $2\ \text{Msps}$ (a $2\ \text{MHz}$
processing bandwidth) that a modest processor can demodulate in real time.
\end{workedexample}

\begin{keytakeaways}
\begin{itemize}[leftmargin=1.4em,itemsep=2pt]
\item SDR digitizes wide and does filtering/mixing/demod in software --- one radio, many waveforms.
\item Front ends: direct RF sampling (fast ADC, bandpass sampling) or direct conversion (I/Q to baseband).
\item A digital downconverter (NCO + mixer + decimating filter) selects and narrows the channel.
\end{itemize}
\end{keytakeaways}

\begin{exercises}
\item Raw data rate of a $61.44\ \text{Msps}$, $12$-bit, complex ADC stream?
\item What decimation factor turns $100\ \text{Msps}$ into a $1\ \text{MHz}$ channel?
\item Name one capability SDR gains from doing conversion digitally.
\end{exercises}

\furtherreading{Steer~\cite{steer}; Razavi~\cite{razavi} (ch.~4).}
